\documentclass[12pt,a4paper]{article}

\usepackage[fontset=none]{ctex}
\setCJKmainfont{Source Han Serif SC}

\usepackage{geometry}
\usepackage{setspace}

\usepackage{booktabs}
\usepackage{enumitem}
\usepackage{hyperref}
\usepackage{datetime}

\usepackage[backend=biber]{biblatex}
\addbibresource{references.bib}


\geometry{left=2.5cm,right=2.5cm,top=2.5cm,bottom=2.5cm}

\onehalfspacing

\hypersetup{
    colorlinks=true,
    linkcolor=black,
    citecolor=blue,
    urlcolor=blue,
    unicode=true
}


\begin{document}

\begin{center}
    \Large \textbf{具有欺骗性的智力量表？} \\[0.5em]
    \Large \textbf{Deceptive Intelligence Scales?} \\
    \vspace{2em}
    \normalsize 潘陶然（Yptrsd） \\[0.5em]
    \normalsize \number\year 年 \number\month 月 \number\day 日 \\[0.5em]
    \normalsize V1.0
\end{center}

% 目录
\tableofcontents

\vspace{1.5em}
\centerline{\hrulefill}
\vspace{1.5em}

% 正文
\paragraph{前言}瑞文量表被认为是用来推断智力水平的一种标准。但是，瑞文量表可能具有潜在的干扰因素，因此它或许并不是衡量智力的唯一标准。瑞文量表在某种程度上疑具有“欺骗性”（这里指其他因素对其的干扰），因而，即使量表的设计本身是具有明确的科学依据的，其使用和结果含义的判读也需要我们更加严谨的对待。

\section{一次偶然的经历}

	在2025年8月，我因为学习上遇到了一些障碍和困难前往了精神卫生中心就诊。在与医生面谈之后，我被开具了瑞文标准推理测验的检查。在全长约14分钟的测试过程中，我发现了题目的图形组合之间的一些规律，因而开始尝试使用我得出的规律对后续题目的答案进行推测。在当时的测试中，后续的题目几乎全部符合我所观测到的规律。最终结果令我非常惊讶：在总共60题的测试中我正确回答了其中的57题，经过综合判定给予我了131的智商测定结果。这不禁引起了我的思考：这个结果究竟是因为我真的有这么高的智商还是因为我成功掌握了题目推理的规律？进一步讲，如果我总结出来这些规律并对一个受试者进行这些规律的培训，能不能让这位受试者获得远超于其真正智力程度或者接近满分的测量结果？再进一步而言，单凭这个量表去界定人聪明与否的行为是否有缺陷，是否不严谨？

\section{为何会如此？}

	在《Encyclopedia of Clinical Neuropsychology》（《临床神经心理学百科全书》）\cite{Leavitt2011}中，我们可以了解到关于其测试逻辑的描述确有“要求受试者推断一个规则来产生这个系列的下一个物体”、“判断所呈现物是否符合存在的规律”、“试题难度渐进，依托于其在测试中积累的知识”的字样。显然，这一智力量表确实在考验人推理规律的能力并以此作为判断智力高低的依据。
%The test requires examinees to infer a rule to generate the next items in a series, or to determine whether a presented design is consistent with the rule. Items become progressively more difficult, building upon knowledge accumulated from the test. As such, the test was designed to assess learning as well as inductive reasoning. -> https://link.springer.com/rwe/10.1007/978-0-387-79948-3_1069

	我们不妨假设瑞文标准推理测验是一个考验“学习——应用”过程的测试，测试结果基于人在任务中的表现。很大程度上，量表的结果可能会受受试者发现一系列规律的早晚而产生不同的变化。测试结果较高的受试者可能在早期甚至首个问题就发现题目中潜在的规律并加以正确的运用，最终得到正确的结果并获得较高测试结果；相反，测试结果较低的受试者则没有能在早期识别规律的能力甚至根本无法识别规律，无法得出正确结果。因此，这种方法在可以保证受试者不会随机选到正确的选项的情况下，从理论出发，针对测试结果较低的受试者是完全可行的。
    
    但是，基于规则认识的测试过程反而带来了一定的不准确性。如果令一位智力水平正常的受试者提前被培训相关的规则（诸如同一行列的图形大小、个数、形状关系，异行列的区别和联系等），那么他将有极大的可能获得较高的测试结果，但由于一些不可抗力因素存在，达成满分或许存在一定困难。无论如何，这类量表就此而言的确是易受干扰的，但一般情况下其设计可以评价受试者在智力层面的能力水平。在测试者不知道量表设计原理和其中的规律时，它可以被认为是基本有效的；但如果出现推理总结能力显著超于其他能力维度的受试者，其结果与理论值相比有可能会浮动。

	显然，作为一套通过逻辑推理来测试受试者智力商数的题目系统，一旦掌握其核心逻辑，就几乎相当于知道了所有题目的答案。类似地，公务员、美术工作者、理学家、计算学家、逻辑学家、哲学家等一系列对图形或逻辑高度敏感的人群由于职业特点和职业的筛选效应，可能会得到相比于其他职业的受试者更高的分数。由于筛选效应，智商更高的人进入这些行业的相对概率也会更高，而非其职业塑造了受试者的职业特点。关于职业特点和筛选效应分别的影响有待进一步探究。
    
\section{验证猜测：可能的实验思路和缺陷}

	为了验证提前预知题目逻辑或部分职业特点对测试结果的影响，我们不妨随机召集一群足够数量以至于可以忽略某些特定影响的未进行过类似测试的志愿者，记录其职业并进行如下分组处理：令每位受试者同时进行瑞文标准推理测试，并将志愿者采用蛇形分配法均分为三组，保证各组基线水平无显著差异。需要注意的是，志愿者的首次分组可以同时使用其他的智商测评工具来减小误差。

    我们需要准备足够数目的相关专家分析量表所示题目的所有解题逻辑，并将其对其中一组受试者进行详细的培训和传达，以所有受试者都能充分理解题目规则为宜。然后，再对其中两组（包括接受培训的一组和其他两组中的任意一组）的志愿者再次进行量表测试，并将最终结果汇总比对。

    经过以上的步骤，我们成功得到了三组志愿者的数据：一组是只接受过一次量表测试的（记为对照组），另一组是接受过两次量表测试但未接受过培训的（记为实验组A），最后一组是接受过两次测试且接受过培训的（记为实验组B）。由于相关流程存在，我们认为三组的第一次测试成绩的基线水平相当。理论上，测试组B的结果相较于对照组和实验组A将会出现显著提升，同时实验组A的结果较对照组也可能出现一定量的上浮（记为结果1）。同时，上文提到的相关具有特殊职业特点的志愿者首次测试的成绩的均分有一定可能会高于其他职业的志愿者，实验组A中具有职业特点人群，其两次测试的涨幅也会相较其他职业更大（记为结果2）。

    \paragraph{结果1}由于经过特别的培训，实验组B可以了解到题目的通用解题方法并得到正确答案。同时，实验组A中的部分志愿者可能在第一次测试的中途成功总结出规律但未能应用到全卷，在其第二次测试时应用到了全卷并获得了正确的结果导致结果上浮。这将证明关于提前预知题目逻辑对测试结果存在影响。
    
    \paragraph{结果2}实验中具有特殊职业特质的人群对于图形和逻辑的敏感度导致其测试结果上浮，也可能因其学习能力更强导致全面了解规则的概率更高，最终在没有培训的情况下在第二次测试中获得更高的测试结果概率大大加大。这将证明职业特质对测试结果存在影响。

    \paragraph{存在的缺陷}整个志愿者群体很难做到像生物实验一般将其他无关变量完全控制，且整个答题过程中存在诸多不稳定、不可预料的因素（诸如紧张等），志愿者本身的状态也对结果有很大影响。因此，辅以其他的手段进行其他的检测是非常必要的，但由于研究水平有限，此类方法的使用还有待细究和研判。

\section{展望}

	据以上来看，直接将瑞文量表的结果不加说明地认为是一个人的智力的行为值得警惕。瑞文量表的结果不直接等同于聪明与否，但是它是受试者整体的智力功能在推理这一角度的投影，并且存在一些不稳定的因素。这种依赖单一因素的测评工具有可能因培训、重复测试等因素导致结果的误差。
    
	就目前而言，若想彻底规避上述提到的问题，我们可能需要一个原理相异的智力评测系统来排除上述的一些干扰。这种系统可能将不再基于量表、题目或者一切依赖受试者反应、操作的流程，而是通过直接对受试者大脑的运作或其他层面进行理化分析来得到结果。但是，这种新系统也可能存在其他尚不明确的问题，因此神经科学、脑科学和心理学的充分研究将是这套新系统的必要前提。若想对现有的测试能力进行一定改进，可以同时加入脑电、眼动等其他辅助手段，而这也有待后续研究。

% 参考文献（可选）
\printbibliography[title={参考文献}]

\paragraph{后记} 把智力量表结果与智力进行简单的等同是一个及其容易掉进去的认知习惯，笔者也曾经陷入过这样的误区。这篇文章初具雏形时，其本意是从本人的亲身经历出发思考瑞文量表用来测试智商的严谨性。在文章写作中期，在进行了反思和讨论之后，笔者终于意识到了这个认识本身即存在一定的误差，于是经历几版修订后遂成文。无论如何，这类量表仅仅是一类工具，其没有绝对的好坏之分，其最终的效用取决于其被利用的方式。因此，一张量表的结果不能决定一个人聪明与否，人类的思维具有显著多样性，每个人在不同的领域都有不同的可能。

\end{document}